% 076_RSAtest_De_ch.tex -- Kapitelinhalt Dok. 076
% Empirische Tests von Periodenfindungsverfahren
% Vollständige Neufassung August 2026 (R75)

\section{Empirische Tests von Periodenfindungsverfahren}

\subsection*{Zusammenfassung}

Dieses Dokument beschreibt eine Testreihe zur Ganzzahlfaktorisierung über
Periodenfindung. Getestet werden klassische Verfahren (Probeteilung,
Pollard~$\rho$) gegen resonanzbasierte Suchverfahren mit unterschiedlichen
Rasterweiten. Das Ergebnis ordnet die Rasterparameter korrekt ein: Sie sind
numerische Abtastgrößen, keine physikalischen Resonanzparameter.

\begin{keyresult}[Getrennte Parameter]
Die Rasterweite der Suchverfahren wird hier mit $\sigma$ bezeichnet
(Werte $1/42$, $1/100$, $1/1000$, $1/100000$). Sie ist \emph{nicht} der
Geometrieparameter $\xi = 4/30000$ der FFGFT. Frühere Fassungen verwendeten
für beide dasselbe Symbol; die Trennung ist notwendig, weil die Größen nichts
miteinander zu tun haben.
\end{keyresult}

\section{Testaufbau}

\subsection{Getestete Verfahren}

\begin{center}
\begin{tabular}{lll}
\toprule
Bezeichnung & Prinzip & Parameter \\
\midrule
\texttt{trial\_division} & Probeteilung bis $\sqrt{N}$ & --- \\
\texttt{pollard\_rho} & Zyklensuche (Floyd) & --- \\
\texttt{period\_classic} & Sequenzielle Periodensuche & $\sigma = 1/100000$ \\
\texttt{period\_universal} & Sequenzielle Periodensuche & $\sigma = 1/100$ \\
\texttt{period\_medium} & Sequenzielle Periodensuche & $\sigma = 1/1000$ \\
\texttt{period\_special} & Sequenzielle Periodensuche & $\sigma = 1/42$ \\
\bottomrule
\end{tabular}
\end{center}

Die Rasterweite $\sigma$ bestimmt die Schrittweite, in der der Suchraum
durchlaufen wird. Sie hat keinen Bezug zu physikalischen Eigenfrequenzen.

\subsection{Testkategorien}

\begin{enumerate}
  \item \textbf{Kleine Semiprimzahlen} ($N < 1000$): Vollständige
    Faktorisierung erwartet, Laufzeitvergleich.
  \item \textbf{Mittlere Semiprimzahlen} ($10^3 < N < 10^6$):
    Skalierungsverhalten.
  \item \textbf{Zwillingsprimprodukte} ($N = p(p+2)$): Verhalten bei
    strukturell benachbarten Faktoren.
  \item \textbf{Große Semiprimzahlen} ($N > 10^6$): Erwarteter Ausfall der
    Periodensuche.
\end{enumerate}

\section{Ergebnisse}

\subsection{Kleine Semiprimzahlen}

Alle Verfahren faktorisieren vollständig. Die Probeteilung ist bei diesen
Größen am schnellsten; die Periodensuche trägt hier keinen Vorteil, da der
Suchaufwand für $r$ den Aufwand der direkten Teilung übersteigt.

\subsection{Skalierungsverhalten}

\textbf{[K]} Mit wachsendem $N$ fällt die Erfolgsrate der Periodensuche ab.
Der Grund ist die Zahl der zu unterscheidenden Kandidaten: Die Periode $r$
teilt $\lambda(N) = \mathrm{lcm}(p-1,q-1)$ und kann Werte bis zur Größenordnung
$N$ annehmen. Eine sequenzielle Suche benötigt damit $O(N)$ Schritte -- deutlich
schlechter als die Probeteilung mit $O(\sqrt{N})$.

\subsection{Einfluss der Rasterweite}

\textbf{[K]} Unterschiedliche $\sigma$-Werte liefern unterschiedliche
Erfolgsraten. Diese Unterschiede sind \emph{Abtastartefakte}: Ein feineres
Raster findet mehr Kandidaten, kostet aber mehr Schritte. Es gibt kein
\enquote{optimales} $\sigma$ im physikalischen Sinn -- nur einen Kompromiss
zwischen Auflösung und Aufwand.

\begin{keyresult}[Warum $\sigma$ keine Resonanzen erzeugt]
Der Bikohärenztest an den Riemann-Nullstellen (Dok.~316, August 2026) zeigt:
Echte Phasenkopplung tritt ausschließlich bei Primzahlpotenzen auf
($b^2 = 0{,}76$--$0{,}85$), nicht bei zusammengesetzten Zahlen
($b^2 = 0{,}018$--$0{,}025$). Alle verwendeten $\sigma$-Werte haben
zusammengesetzte Nenner ($42 = 2\cdot3\cdot7$, $100 = 2^2\cdot5^2$,
$1000 = 2^3\cdot5^3$). Sie bilden im Spektrum \emph{Mischpunkte} ohne eigene
Kopplungslinie. Ihre Wirkung ist rein numerisch.
\end{keyresult}

\subsection{Zwillingsprimprodukte}

Bei $N = p(p+2)$ liegen die Faktoren dicht beieinander. Die Probeteilung
profitiert davon (Fund nahe $\sqrt{N}$), die Periodensuche nicht -- die
Periode $r$ hängt von $\mathrm{lcm}(p-1,p+1)$ ab und zeigt kein besonderes
Muster. Die in früheren Fassungen beschriebene \enquote{Zwillingsprim-Optimierung}
über eine spezielle Rasterweite ist nicht reproduzierbar.

\section{Einordnung}

\subsection{Algorithmenvergleich}

\begin{center}
\begin{tabular}{llll}
\toprule
Verfahren & Prinzip & Schließungsbedingung & Komplexität \\
\midrule
Probeteilung & Direkte Division & Diskrete Arithmetik & $O(\sqrt{N})$ \\
Pollard $\rho$ & Zyklensuche & Geburtstagsparadoxon & $O(N^{1/4})$ \\
Periodensuche & Spektralsuche & Rationalisiertes Raster & $O(N)$, Weyl-Grenze \\
Shors QFT & Quanteninterferenz & Diskrete Phase & $O((\log N)^3)$ auf QC \\
\bottomrule
\end{tabular}
\end{center}

\textbf{[K]} Die klassische Periodensuche ist den etablierten Verfahren
unterlegen. Sie hat keinen Komplexitätsvorteil; die früher behauptete
Übereinstimmung mit Shors $O((\log N)^3)$ gilt allein für den Quantenfall und
ist für ein klassisches Verfahren nicht belegt.

\subsection{Die strukturelle Grenze}

\textbf{[K]} Über die praktische Unterlegenheit hinaus besteht eine
prinzipielle Grenze. Nach der Weyl-Obstruktion (Dok.~316) wächst das Spektrum
jedes kompakten Resonanzraums wie $N(\lambda) \sim \lambda^{d/2}$, während die
arithmetische Zähldichte wie $T\ln T$ wächst. Ein endlicher Resonanzraum kann
diese Dichte nicht tragen -- unabhängig von Rasterweite und Rechengenauigkeit.

\subsection{Primzahlen als Generatoren}

\textbf{[K]} Was die Resonanzstruktur tatsächlich trägt, sind die Primzahlen.
Das Euler-Produkt behandelt jede Primzahl für sich, ohne Mischterme; das duale
Längenspektrum enthält nur $k \ln p$. Die Naturtonreihe zeigt dieselbe Struktur
physikalisch: In der Reihe $n f_0$ führen genau die Primzahl-Ordnungen neue
Klangqualitäten ein, $4 = 2\cdot2$ und $6 = 2\cdot3$ sind Kombinationen.

Ein Verfahren, das Primzahlen als orthogonale Achsen behandelt, arbeitet
strukturell korrekt. Ein Verfahren, das eine zusammengesetzte Rasterweite als
Resonanzparameter behandelt, tut das nicht.

\section{Bilanz}

\begin{center}
\begin{tabular}{ll}
\toprule
Aussage & Status \\
\midrule
Periodenfindung ist ein Spektralproblem & \textbf{[K]} \\
Primzahlen tragen die Resonanzstruktur & \textbf{[K]} bikohärenzgemessen \\
Aufbereitung macht ~95\,\% der Gatterkosten aus & \textbf{[K]} (Dok.~176) \\
Sequenzielle Periodensuche ist $O(N)$ & \textbf{[K]} gemessen \\
$\sigma$-Werte sind Abtastraster & \textbf{[K]} \\
$\sigma$-Werte erzeugen Eigenfrequenzen & \textbf{[X]} zusammengesetzt \\
Klassische Periodensuche erreicht $O((\log N)^3)$ & \textbf{[X]} nicht belegt \\
Zwillingsprim-Optimierung über Rasterweite & \textbf{[X]} nicht reproduzierbar \\
Skalierung durch Rechengenauigkeit lösbar & \textbf{[X]} Weyl-Obstruktion \\
\bottomrule
\end{tabular}
\end{center}

\textbf{Der tragende Satz.} Die Testreihe dokumentiert das Verhalten
sequenzieller Periodensuchverfahren. Sie zeigt keinen Vorteil gegenüber
etablierten Methoden, und die beobachteten Unterschiede zwischen Rasterweiten
sind Abtastartefakte. Die Resonanzbeschreibung selbst bleibt richtig -- nur
sind ihre Träger die Primzahlen, nicht das Raster.

\section{Bezug zum Korpus}

Die Einordnung stützt sich auf Dok.~316 (August 2026): Bikohärenztest,
Weyl-Obstruktion und die Rolle der Primzahlen als Generatoren. Die
Überarbeitung ist unter R75 in Dok.~190 deklariert.

\section{Verifikation}

Reproduzierbar durch \texttt{s1\_periodensuche.py} (Korrektheit,
Skalierung, $\sigma$-Zerlegung, Verfahrensvergleich) und
\texttt{s2\_grenzen.py} (Grenzen beider Ansätze). Reine
Standardbibliothek, Sollwerte als Assertions, Seed 20780458.
